\documentclass[11pt]{article}
\newcommand{\documenttype}{Research statement}
\usepackage[letterpaper,margin=1in]{geometry}
\usepackage[T1]{fontenc}
\usepackage[utf8]{inputenc}
\usepackage{newtxtext,newtxmath}
\usepackage{microtype}
\usepackage{xcolor}
\usepackage{enumitem}
\usepackage{tabularx}
\usepackage{titlesec}
\usepackage{fancyhdr}
\usepackage{hyperref}
\definecolor{accent}{HTML}{315B8A}
\definecolor{link}{HTML}{5B4F7A}
\hypersetup{colorlinks=true,urlcolor=link,linkcolor=link,citecolor=link,
  pdfauthor={Phurailatpam Hemantakumar}}
\setlength{\parindent}{0pt}
\setlength{\parskip}{6pt}
\setlength{\headheight}{14pt}
\setlength{\emergencystretch}{1.5em}
\titleformat{\section}{\large\bfseries\color{accent}}{}{0pt}{}
\titlespacing*{\section}{0pt}{11pt}{5pt}
\titleformat{\subsection}{\normalsize\bfseries}{}{0pt}{}
\titlespacing*{\subsection}{0pt}{8pt}{3pt}
\setlist[itemize]{leftmargin=15pt,itemsep=3pt,topsep=3pt,parsep=0pt}
\setlist[enumerate]{leftmargin=19pt,itemsep=7pt,topsep=4pt,parsep=0pt}
\pagestyle{fancy}
\fancyhf{}
\fancyhead[L]{\small \href{https://hemantaph.com}{Phurailatpam Hemantakumar}}
\fancyhead[R]{\small\documenttype}
\fancyfoot[R]{\small\thepage}
\renewcommand{\headrulewidth}{0.3pt}
\renewcommand{\footrulewidth}{0pt}
\newcommand{\firstheading}[2]{%
  \thispagestyle{plain}%
  {\LARGE\bfseries\color{accent} #1}\par
  \vspace{3pt}{#2}\par
  \vspace{4pt}{\color{accent}\hrule}\vspace{8pt}}
\newcommand{\dated}[2]{\par\noindent\begin{tabularx}{\linewidth}{@{}X r@{}}#1 & #2\end{tabularx}\par}

\setlength{\parskip}{3pt}
\titlespacing*{\section}{0pt}{9pt}{4pt}
\titlespacing*{\subsection}{0pt}{7pt}{3pt}
\hypersetup{pdftitle={Research statement | Phurailatpam Hemantakumar}}
\begin{document}
\firstheading{Research statement}{\href{https://hemantaph.com}{Phurailatpam Hemantakumar}}

Strong lensing can magnify distant gravitational-wave (GW) sources and produce repeated signals. My lensing work estimates detectable populations and derives population priors for candidate assessment.

Separately, I co-develop and implement an analytic, frequency-domain post-Newtonian (PN) eccentric waveform model with collaborators, for eccentricity inference and studies of compact-binary formation.

My multimessenger work forecasts the detectability of lensed gamma-ray bursts and GW counterparts. I am developing statistical tests to distinguish lensed repetitions from unrelated bursts.

\section*{Research foundations and ongoing work}
\subsection*{Gravitational-wave lensing statistics}
I lead development of \href{https://ler.hemantaph.com/}{\texttt{ler}}, an open-source framework for simulating lensed and unlensed compact-binary populations and detectable rates [\hyperlink{ref1}{1}], with Otto Akseli Hannuksela, Anupreeta More, David Keitel, Chris Van Den Broeck, and other collaborators. My work combines source and lens populations, strong-lensing selection, lens-equation solutions, image properties, and detector selection. For elliptical power-law lenses with external shear, I implemented cross-section and optical-depth calculations and semi-analytic lens-equation solutions. These calculations support predictions of the magnification and time-delay distributions relevant to lensing searches.

To reduce the cost of repeated detector-selection calculations, I developed \href{https://gwsnr.hemantaph.com/}{\texttt{gwsnr}} with Otto Akseli Hannuksela [\hyperlink{ref2}{2}]. It provides partial-scaling interpolation and parallel waveform inner products for signal-to-noise ratios, together with neural-network detectability estimates. These tools have supported collaborative lensing follow-up, faint-counterpart searches, and population terms in Bayesian lensing posterior odds [\hyperlink{ref3}{3},\hyperlink{ref4}{4}]. I have implemented propagation of population-hyperparameter and finite-sampling uncertainties in \texttt{ler}; further analysis is ongoing.

\subsection*{Eccentric waveform development and implementation}
Building on Tiwari et al. (2019) [\hyperlink{ref10}{10}], which I co-authored, I co-developed \href{https://arxiv.org/abs/2508.12697}{\texttt{TaylorF2Ecck}} with Achamveedu Gopakumar, Maria Haney, Tjonnie Li, and Srishti Tiwari. I implemented the model in a \texttt{LALSuite} fork [\hyperlink{ref5}{5}]. This nonspinning, low-eccentricity inspiral model includes leading eccentricity contributions from orbital dynamics, periastron advance, and GW emission, with phasing through 3PN order. I conducted Bayesian analyses of GW170817 and GW190425 and injection studies. We found no strong evidence favouring eccentricity or measurable periastron-advance effects. Our accuracy studies identified potential mass-ratio biases from insufficient phasing accuracy, motivating eccentric phasing through at least 3.5PN order. In a separate project with Lalit Pathak and Achamveedu Gopakumar, I studied eccentricity and line-of-sight acceleration in GW190814. We found a strong degeneracy between the two effects, without significant evidence for either~[\hyperlink{ref6}{6}].

\subsection*{Multimessenger lensing}
With Anupreeta More, I developed detectability forecasts for lensed short gamma-ray bursts (GRBs) and binary-neutron-star GW counterparts using \texttt{ler} and \texttt{gwsnr} [\hyperlink{ref7}{7}]. This work examined both GW counterparts to lensed GRBs and GRB counterparts to lensed GW systems. In collaboration with Otto Akseli Hannuksela and \href{https://iphd.site.nthu.edu.tw/p/406-1222-87294,r5205.php?Lang=en}{Albert Kong}, I am now developing Bayesian tests of whether multiple GRBs are lensed images of one source, initially using sky localisation and time delays. This adapts the GW lensing posterior-odds framework [\hyperlink{ref4}{4}], with GRB-specific uncertainties and selection.

\newpage
\section*{Proposed research}
\subsection*{1. Lensing populations, candidate assessment, and cosmography}
\textbf{Scientific objective.} I aim to determine how reliably we can predict detectable lensed GW populations and use them in astrophysical inference. Source evolution, lens populations, and selection all influence the observed sample. Their uncertainties affect both the expected event rate and the prior information used to assess whether candidate signals are lensed.

\textbf{Merger histories and lensing statistics.} With \href{https://pure.ewha.ac.kr/en/persons/chunglee-kim/}{Chunglee Kim}, I am working on binary-black-hole merger histories informed by star formation, metallicity-dependent formation efficiency, and formation-to-merger delays. I will extend this work by convolving metallicity-weighted star formation with delay-time distributions. I will constrain these models with LVK observations while including detector selection, and compare alternative histories rather than treating the high-redshift merger rate as uniquely determined by nearby detections. Within \texttt{ler}, I will propagate source and lens hyperparameter uncertainties and compare population-model choices. I will distinguish these astrophysical uncertainties from finite observed-catalogue statistics and finite Monte Carlo error.

The first outputs will be detectable rates and magnification and time-delay distributions with quantified uncertainties. I will use these to test the sensitivity of Bayesian lensing posterior odds to rate and population priors. For detector selection, I will benchmark rapid \texttt{gwsnr} estimates against direct calculations near the detection boundary and compare SNR-based probabilities with injection recovery. These comparisons will assess how well the fast estimates describe search sensitivity and where calibration is needed. \textbf{LVK/GRASP commitments.} I will continue project leadership, development, testing, documentation, and user support for \texttt{ler}, including ISR-weighted work towards GWTC-6 and GWTC-7.

\textbf{Extensions of the current \texttt{ler} statistical work.} In collaboration with Neha Singh, I will provide simulated field-binary populations for rate estimates and study the signatures of lensed binary-black-hole mergers originating from Population III stars. With Anupreeta More, I will extend the lensing-statistics calculations to galaxy-cluster lenses.

\textbf{Population-based time-delay cosmography.} In a related proposed project, I will test the robustness of cosmological inference from the lensed-event fraction and time-delay distribution, building on Maity et al. [\hyperlink{ref8}{8}]. I will recompute these observables with alternative source and lens populations, elliptical power-law lenses with external shear, and detector selection. Simulated catalogues will test parameter bias, credible-interval coverage, and degeneracies between cosmology and population assumptions. The aim is to identify which constraints remain reliable when astrophysical uncertainties are included and which require stronger external population information. This is a forecast and validation programme; its initial outcomes do not depend on a confirmed lensed detection.

\subsection*{2. Eccentric waveform development and implementation}
\textbf{Scientific objective.} My goal is to extend and implement our analytic, frequency-domain PN eccentric waveform model, \texttt{TaylorF2Ecck}, with collaborators. The model can be used to infer orbital eccentricity and investigate compact-binary formation channels.

Further development of \texttt{TaylorF2Ecck} is a collaborative effort. Shri Krishna Anand, with guidance from Achamveedu Gopakumar, is incorporating spin effects into our eccentric waveform prescription, targeting 3PN phasing through $\mathcal{O}(e_0^6)$ and 1PN amplitudes through $\mathcal{O}(e_0^4)$. My planned contribution is to implement these developments in a \texttt{LALSuite} fork and test and validate them through analytical checks, comparisons with independent models, and injection studies. I plan to include eccentric waveform development and implementation in my GRASP activities and pursue integration into the official \texttt{LALSuite} repository.

\newpage
\textbf{Higher-order eccentric phasing.} Our published accuracy studies motivate eccentric phasing through at least 3.5PN order. Chandra Kant Mishra and Kaushik Paul may also be developing such corrections. I will seek coordination with them; otherwise, I will undertake the development with external collaborators. I plan to incorporate these corrections into \texttt{TaylorF2Ecck}.

\subsection*{Waveform extensions: line-of-sight acceleration}
Lalit Pathak, Achamveedu Gopakumar, and I plan to extend our GW190814 analysis to more events, using higher-order eccentric corrections and a more general treatment of time-dependent source acceleration, such as Zhao et al.~[\hyperlink{ref9}{9}]. Joint inference and injection studies will assess how well eccentricity and acceleration can be distinguished and how waveform assumptions affect the results.

\subsection*{3. Multimessenger lensing and GRB candidate tests}
\textbf{Scientific objective.} As a follow-up to More and Phurailatpam (2025) [\hyperlink{ref7}{7}], I aim to extend our detectability studies to quantitative tests of whether repeated GRBs are images of the same lensed source. Such tests can guide candidate follow-up and, when a GW counterpart is possible, targeted searches. The immediate statistical problem is to distinguish lensing from unrelated bursts with apparently compatible properties.

I will complete and validate the sky-localisation and time-delay posterior-odds calculation following Hannuksela et al. (2026) [\hyperlink{ref4}{4}], including measurement uncertainties, unrelated-burst coincidences, lensing-rate priors, and GRB selection. Simulated lensed and unrelated populations will test candidate ranking and false-association rates. I will then incorporate light-curve information and explore joint GW--GRB associations, building on my GRB--GW collaboration with Otto Akseli Hannuksela and \href{https://iphd.site.nthu.edu.tw/p/406-1222-87294,r5205.php?Lang=en}{Albert Kong}. I will check how each additional observable changes the inference. Our published detectability forecasts, my ongoing implementation, and experience with lensing population priors provide the basis for this work.

The near-term goal is to test the statistical framework and release documented code for candidate assessment. Longer-term applications to source and lens properties, or time-delay cosmography of an individually identified system, would require suitable electromagnetic observations and lens modelling. These applications are distinct from the GW population-based cosmography study above.

\section*{Feasibility and research programme}
These projects build on existing software, published analyses, and established collaborations. The planned outcomes are lensing forecasts with quantified population uncertainties, implemented and validated eccentric waveform extensions, and a tested GRB lensing framework. I will lead the population-simulation work and contribute waveform implementation and Bayesian analyses within the respective collaborations, releasing documented code and reproducible examples.

\section*{References}
\begingroup
\setlength{\parskip}{0pt}
\small
\hypertarget{ref1}{[1]} H. Phurailatpam et al. (2026), \textit{ler}, JOSS \textbf{11}, 7420. \href{https://doi.org/10.21105/joss.07420}{doi:10.21105/joss.07420}.\par
\hypertarget{ref2}{[2]} H. Phurailatpam and O. A. Hannuksela, \textit{gwsnr}, \href{https://arxiv.org/abs/2412.09888}{arXiv:2412.09888}.\par
\hypertarget{ref3}{[3]} L. C. Y. Ng et al. (2025), faint lensed GW counterparts, MNRAS \textbf{540}, 2937. \href{https://doi.org/10.1093/mnras/staf874}{doi:10.1093/mnras/staf874}.\par
\hypertarget{ref4}{[4]} O. A. Hannuksela et al. (2026), lensing posterior odds, ApJ \textbf{1002}, 42. \href{https://doi.org/10.3847/1538-4357/ae5816}{doi:10.3847/1538-4357/ae5816}.\par
\hypertarget{ref5}{[5]} H. Phurailatpam et al. (2026), eccentric inspiral inference, PRD \textbf{114}, 063003. \href{https://doi.org/10.1103/94fx-wnjx}{doi:10.1103/94fx-wnjx}.\par
\hypertarget{ref6}{[6]} L. Pathak, H. Phurailatpam, and A. Gopakumar (2026), tertiary object in GW190814, \href{https://arxiv.org/abs/2605.21955}{arXiv:2605.21955}.\par
\hypertarget{ref7}{[7]} A. More and H. Phurailatpam (2025), multimessenger lensing, Phil. Trans. R. Soc. A \textbf{383}, 20240127. \href{https://doi.org/10.1098/rsta.2024.0127}{doi:10.1098/rsta.2024.0127}.\par
\hypertarget{ref8}{[8]} K. N. Maity et al., strong-lensing cosmography with binary black holes, \href{https://arxiv.org/abs/2512.15168}{arXiv:2512.15168}.\par
\hypertarget{ref9}{[9]} X. Zhao, H. Yan, and X. Chen (2026), waveforms for accelerating sources, PRD \textbf{114}, 063006. \href{https://doi.org/10.1103/swkb-nkgg}{doi:10.1103/swkb-nkgg}.\par
\hypertarget{ref10}{[10]} S. Tiwari et al. (2019), eccentric inspiral templates, PRD \textbf{99}, 124008. \href{https://arxiv.org/abs/1905.07956}{arXiv:1905.07956}.\par
\endgroup
\end{document}
